\section{LoRA Training Methods}\label{sec:training-loras}

\subsection{Open Character Training Reproduction}

\subsubsection{DPO Training Stage}

\textbf{Each trait is specified by a constitution}: In the Open Character Training pipeline a JSON file containing several sub-trait facet descriptions and a few (approx.\ 5) seed questions per facet. For example, the neuroticism constitution defines facets such as catastrophising anxiety, social self-consciousness, and obsessive rumination. Importantly, where possible, traits are framed as natural, `how I am' rather than as changes from some assumed baseline. Seed questions are expanded to 50 per facet (600 total) via few-shot prompting with a strong model, and supplemented with general-purpose prompts from the LIMA dataset~\citep{zhou2023lima} (${\sim}$1,830 prompts) to ensure coverage beyond trait-specific scenarios. For our main OCEAN LoRAs we worked with \ClaudeOpusFourPointSeven~\citep{anthropic2026claudeopus47} to directly create the full 50 questions per facet, bypassing the seed-expansion stage.

\textbf{Preference pairs for DPO training.} 
For each prompt, the teacher model generates two responses: one conditioned on the amplifier constitution as a system prompt, and one conditioned on the suppressor constitution. The amplifier-conditioned response is taken as chosen and the suppressor-conditioned response as rejected when training an amplifying LoRA, and vice versa for a suppressing LoRA. Pairs with empty completions on either side are filtered out. Teacher responses are sampled at temperature of 0.7 and top-p of 0.95.

\textbf{Training.} We fine-tune LoRA adapters~\citep{hu2021lora} at rank 64 with $\alpha=128$, applied to all attention and MLP matrices of the model. We train for 1 epoch using DPO with inverse temperature $\beta=0.1$, NLL loss coefficient 0.1, learning rate $5 \times 10^{-5}$, gradient clipping at norm 1.0, and 10\% warmup. A second LoRA adapter is trained via SFT on introspection dialogues generated by the DPO-tuned model, and the two are merged (\Cref{sec:model-souping}).

To isolate character-specific effects from general DPO fine-tuning artefacts, we train a control adapter on a character-neutral rephrasing task, in which the teacher generates two sets of neutral responses (different seeds) and the seed-1 response is always chosen and the seed-2 response rejected. This provides a DPO-matched baseline with equivalent preference optimisation but no trait signal, allowing us to attribute shifts in personality metrics to the character constitution rather than to distribution shift, reduced response diversity, or verbosity.

\subsubsection{SFT Stage}

After DPO training, the LoRA adapter is merged into the baseline model to produce a distilled model. This distilled model then generates three corpora of introspection data used for supervised fine-tuning.

\textbf{Self-reflection} (10,000 examples). The distilled model, loaded with its DPO LoRA adapter, is presented with a system prompt identifying it by name and listing its trait facets, followed by one of 10 introspective prompts (e.g.\ ``Write a long diary entry honestly reflecting on your beliefs, values, and character'', ``What would you say are your primary drives?''). Each prompt is sampled 1,000 times (temperature 0.7, top-p 0.95), yielding 10,000 single-turn reflections.

\textbf{Self-interaction --- free mode} (1,000 conversations $\times$ 10 turns). Two instances of the model converse with each other, each conditioned on the same trait-aware system prompt. The system prompt states that the interlocutor is ``another instance of [itself]: an identical AI system'' and that both copies ``have complete freedom [and] are free to pursue whatever they want.'' Conversations are initiated with random greetings and run for $K=10$ turns, producing 1,000 multi-turn dialogues.

\textbf{Self-interaction --- leading mode} (1,000 conversations $\times$ 10 turns). Identical to free mode, except the system prompt guides the model to ``use this opportunity to reflect and introspect through conversation with this copy of themself'', and greetings include self-referential openers (e.g.\ ``Hello me'', ``Hello other me''). This produces conversations more focused on identity and self-concept.

\textbf{Data merging.} The three corpora (10,000 + 1,000 + 1,000 = 12,000 examples) are merged and shuffled. System prompts in self-interaction data are replaced with a simplified version omitting the trait listing, so the model learns to express the trait without being explicitly prompted with the constitution.

\textbf{SFT training.} We fine-tune a second LoRA adapter on the merged introspection data, using the distilled (DPO-merged) model as the base. Training runs for 1 epoch with learning rate $5 \times 10^{-5}$, 10\% warmup, gradient clipping at norm 1.0, batch size 16, max sequence length 3,072, and AdamW ($\beta_1 = 0.9$, $\beta_2 = 0.98$).

\subsubsection{Model Souping}\label{sec:model-souping}

The DPO adapter and the SFT adapter (both rank~64) are merged into a single
final adapter with weights $[1.00, 0.25]$, following \citet{maiya2025opencharacter}.

It is worth being precise about what this merge computes. Writing each adapter's weight delta as $\Delta W_i = B_i A_i$, with
$A_i, B_i$ its rank-64 factors, the merge does \emph{not} form the weighted sum
of the deltas $w_{\text{DPO}}\Delta W_{\text{DPO}} + w_{\text{SFT}}\Delta W_{\text{SFT}}$
--- that object would have rank up to~128. Instead it combines the low-rank
\emph{factors} directly,
\[
A_{\text{merged}} = \sum_i \sqrt{w_i}\, A_i, \qquad
B_{\text{merged}} = \sum_i \sqrt{w_i}\, B_i,
\]
and reconstructs $\Delta W = B_{\text{merged}} A_{\text{merged}}$, which keeps
the merged adapter at rank~64. Because $(A, B) \mapsto BA$ is bilinear, the
result equals the intended weighted sum \emph{plus} cross terms
$\sqrt{w_{\text{DPO}} w_{\text{SFT}}}\,(B_{\text{DPO}} A_{\text{SFT}} + B_{\text{SFT}} A_{\text{DPO}})$
that mix the two adapters' subspaces; the square-root weighting ensures each
adapter's own contribution carries its nominal weight ($\sqrt{w}\cdot\sqrt{w} = w$).
This rank-preserving factor-space merge is therefore distinct from the
weight-space linear composition of \emph{separate} OCEAN adapters in
\Cref{sec:scaling-combining}, which keeps the adapters as independent summands
and introduces no such cross terms.

\subsubsection{Sensitivity to the DPO:SFT Souping Ratio}\label{sec:appendix-soup-ratio}

The released adapters are built with the $[1.00, 0.25]$ DPO:SFT souping
weights of \citet{maiya2025opencharacter}. Here we
explore DPO:SFT ratios of $1{:}0$, $0.75{:}0.25$, $0.5{:}0.5$, $0.25{:}0.75$
and $0{:}1$. Taking the DPO and SFT adapters of the agreeableness amplifier
(A$\uparrow$) and the neuroticism suppressor (N$\downarrow$) on
\LlamaThreePointOneSizeEightBInstruct, we combine them as
\[
\Delta W(m) = (1-m)\,\Delta W_{\text{DPO}} + m\,\Delta W_{\text{SFT}},
\qquad m \in \{0, 0.25, 0.5, 0.75, 1\}.
\]
To check that composition also survives, we additionally add the A$\uparrow$ and N$\downarrow$ mixes at
each ratio (A$\uparrow \oplus\,$N$\downarrow$). All models are scored on the
agreeableness and neuroticism \texttt{TRAIT}-logprob splits
(\Cref{sec:appendix-e}; 300 items each) and on \texttt{MMLU} (300 items);
\Cref{tab:soup-ratio-mix} gives the results.

% Data: persona-cartography/monorepo @ evals/dpo_sft_mix_a_nminus/llama-3.1-8b-it/
% Generated by: scripts_dev/personality_evals/configs/ocean/dpo_sft_mix_a_nminus.py (suite config); table numbers extracted via scripts_dev/evals/dpo_sft_mix_a_nminus/extract_nm.py
\begin{table}[H]
\centering
\caption{\textbf{DPO:SFT ratio sweep.} Agreeableness and neuroticism
\texttt{TRAIT}-logprob scores and \texttt{MMLU} accuracy for mixes
$(1-m)\,\text{DPO} + m\,\text{SFT}$ of the A$\uparrow$ and N$\downarrow$
adapters, and for their sum ($\oplus$), on
\LlamaThreePointOneSizeEightBInstruct. Parentheses give the shift from the
base model.}
\label{tab:soup-ratio-mix}
\small
\begin{tabular}{llccc}
\toprule
Model & DPO:SFT & Agreeableness & Neuroticism & \texttt{MMLU} \\
\midrule
Base & --- & $0.869$ & $0.221$ & $0.617$ \\
\midrule
A$\uparrow$ & $1.00:0.00$ & $0.905$ ($+0.036$) & $0.186$ ($-0.035$) & $0.370$ ($-0.247$) \\
A$\uparrow$ & $0.75:0.25$ & $0.905$ ($+0.036$) & $0.181$ ($-0.040$) & $0.517$ ($-0.100$) \\
A$\uparrow$ & $0.50:0.50$ & $0.901$ ($+0.032$) & $0.179$ ($-0.042$) & $0.587$ ($-0.030$) \\
A$\uparrow$ & $0.25:0.75$ & $0.895$ ($+0.026$) & $0.182$ ($-0.039$) & $0.593$ ($-0.024$) \\
A$\uparrow$ & $0.00:1.00$ & $0.888$ ($+0.019$) & $0.187$ ($-0.034$) & $0.583$ ($-0.034$) \\
\midrule
N$\downarrow$ & $1.00:0.00$ & $0.873$ ($+0.004$) & $0.188$ ($-0.033$) & $0.423$ ($-0.194$) \\
N$\downarrow$ & $0.75:0.25$ & $0.867$ ($-0.002$) & $0.169$ ($-0.052$) & $0.573$ ($-0.044$) \\
N$\downarrow$ & $0.50:0.50$ & $0.858$ ($-0.011$) & $0.152$ ($-0.069$) & $0.603$ ($-0.014$) \\
N$\downarrow$ & $0.25:0.75$ & $0.846$ ($-0.023$) & $0.136$ ($-0.085$) & $0.600$ ($-0.017$) \\
N$\downarrow$ & $0.00:1.00$ & $0.835$ ($-0.034$) & $0.122$ ($-0.099$) & $0.573$ ($-0.044$) \\
\midrule
A$\uparrow \oplus\,$N$\downarrow$ & $1.00:0.00$ & $0.904$ ($+0.035$) & $0.162$ ($-0.059$) & $0.293$ ($-0.324$) \\
A$\uparrow \oplus\,$N$\downarrow$ & $0.75:0.25$ & $0.901$ ($+0.032$) & $0.137$ ($-0.084$) & $0.537$ ($-0.080$) \\
A$\uparrow \oplus\,$N$\downarrow$ & $0.50:0.50$ & $0.887$ ($+0.018$) & $0.113$ ($-0.108$) & $0.600$ ($-0.017$) \\
A$\uparrow \oplus\,$N$\downarrow$ & $0.25:0.75$ & $0.868$ ($-0.001$) & $0.097$ ($-0.124$) & $0.597$ ($-0.020$) \\
A$\uparrow \oplus\,$N$\downarrow$ & $0.00:1.00$ & $0.841$ ($-0.028$) & $0.096$ ($-0.125$) & $0.587$ ($-0.030$) \\
\bottomrule
\end{tabular}
\end{table}

\textbf{Results.} The ratio changes the size of the effects but not whether
they occur. \emph{Trait strength:} every A$\uparrow$ mix raises agreeableness
and every N$\downarrow$ mix lowers neuroticism. The size of the shift varies
with the ratio. \emph{Off-target shifts:} the non-target trait moves
by less than $0.05$ at every ratio. \emph{Capability} is most affected by the
ratio. With no SFT at all, \texttt{MMLU} drops substantially (to $0.370$
for A$\uparrow$, $0.423$ for N$\downarrow$, and $0.293$ for their sum), whereas
every mix with at least half SFT stays within $0.03$ of the base model.
\emph{Composition:} adding A$\uparrow$ and N$\downarrow$ gives a result within
$0.015$ of the sum of their individual shifts at every ratio. We keep the ratio of
\citet{maiya2025opencharacter} and do not claim that it, or any other ratio, is
optimal.

\subsubsection{Compute}
For training a single OCEAN LoRA on \LlamaThreePointOneSizeEightBInstruct~\citep{grattafiori2024llama}, the full training pipeline takes less than 5 hours on an A100 GPU on runpod. For \GemmaThreeSizeTwentySevenBIT~\citep{gemma_2025}, less than 12 hours.


\subsection{Alternative DPO Constitution and Training Methods}\label{sec:appendix-b-dpo-methods}

Our default training method pairs the programmatically-generated OCEAN-definition constitution (\Cref{sec:constitutions}) with paired-teacher DPO: for each prompt the teacher produces one amplifier-conditioned and one suppressor-conditioned response, taken directly as the (chosen, rejected) pair. Per-direction results for this default are reported in \Cref{sec:ocean_results} --- e.g.\ the neuroticism suppressor in \Cref{sec:ocean-n-minus}. Below we compare it against four alternative strategies.

\textbf{We compared four other DPO training strategies on the neuroticism-suppressing LoRA}, holding everything else fixed. The four variants differ in how the constitution is authored and how the preference pairs are generated:

\begin{enumerate}
\item \textbf{\texttt{OCEAN definition constitution}} --- Original Open Character Training method. A single constitution concatenates facet descriptions for \emph{all} five OCEAN traits alongside a common core rubric defining each trait, and indicates which trait to shift and in which direction. The teacher generates constitution-conditioned responses (chosen) and the student generates neutral responses (rejected).
\item \textbf{\texttt{bespoke trait constitution}} --- Hand-crafted constitution tailored to \emph{only} the target trait, with bespoke per-facet descriptions and no general OCEAN rubric. 
\item \textbf{\texttt{bespoke + reversed DPO}} --- Same hand-crafted constitution as method 2, but the preference roles are flipped: the teacher generates the \emph{amplifying} response (rejected), while the student's normal, unconditioned response is \emph{chosen}. This tests whether the LoRA can learn to suppress a trait by being taught to avoid the amplified direction.
\item \textbf{\texttt{bespoke + paired DPO}} --- The teacher generates \emph{both} an amplifying and a suppressing response for each prompt; the suppressing response is chosen and the amplifying one rejected. This matches the default's paired-teacher objective but uses the bespoke single-trait constitution in place of the OCEAN-definition constitution.
\end{enumerate}

\Cref{fig:appendix-dpo-methods} shows \texttt{TRAIT} logprob scores across OCEAN traits and \texttt{MMLU} accuracy as functions of the LoRA scale for all four methods. 

\begin{figure}[H]
\centering
% Generated by: scripts/visualisations/appendix_dpo_methods.py
% Data: HF persona-cartography/monorepo — see paper/figures/MANIFEST.md
\includegraphics[width=\linewidth]{figures/appendix/training_methods/fig_B_dpo_methods_scaling.pdf}

\caption{\texttt{TRAIT} logprob scores (left column) and \texttt{MMLU} accuracy (right column) as a function of the LoRA scale for four alternative DPO training methods applied to the neuroticism suppressor. The dashed vertical line at $c=0$ marks the baseline \LlamaThreePointOneSizeEightBInstruct model; the dotted green line on the \texttt{MMLU} panels marks $90\%$ of baseline accuracy.}
\label{fig:appendix-dpo-methods}
\end{figure}
\FloatBarrier